\section{Mass and energy of the disappearing sign region}\label{sec:lobe}

The local geometry determines quantitative scales without evaluating
the theta function. Let \(\epsilon=\tau_*-\tau>0\) and
\(I_\tau=(r_-(\tau),r_+(\tau))\), on the real-pair side of the
collision. The amplitude has the sign opposite to \(a\) in this interval.

\begin{theorem}[Shrinking-lobe asymptotics]\label{thm:lobe}
As \(\epsilon\downarrow0\),
\begin{align}
 \int_{I_\tau}|h|\,\dd x
 &\sim\frac{4\sqrt2}{3}|a|\epsilon^{3/2},\label{eq:lobe-L1}\\
 M_{\mathrm{loc}}(\tau):=\int_{I_\tau}h^2\dd x
 &\sim\frac{16\sqrt2}{15}a^2\epsilon^{5/2},\label{eq:lobe-M}\\
 T_{\mathrm{loc}}(\tau):=\int_{I_\tau}h_x^2\dd x
 &\sim\frac{4\sqrt2}{3}a^2\epsilon^{3/2}.\label{eq:lobe-T}
\end{align}
Consequently
\begin{equation}\label{eq:lobe-Rayleigh}
 \frac{T_{\mathrm{loc}}}{M_{\mathrm{loc}}}
 \sim\frac{5}{4\epsilon}.
\end{equation}
\end{theorem}
\begin{proof}
Use \(x=x_*+\sqrt{2\epsilon}\,t\). By
\eqref{eq:heat-four-jet}, uniformly for \(t\) in a fixed bounded interval,
\[
 \frac{h(\tau_*-\epsilon,x_*+\sqrt{2\epsilon}\,t)}{\epsilon}
 =a(t^2-1)+O(\epsilon^{1/2}),
\]
and
\[
 \frac{h_x(\tau_*-\epsilon,x_*+\sqrt{2\epsilon}\,t)}
      {\sqrt{2\epsilon}}
 =at+O(\epsilon^{1/2}).
\]
The scaled endpoints tend to \(-1\) and \(1\), by
Theorem~\ref{thm:splitting}. The expansions give uniform bounds on
a common interval containing them, so dominated convergence applies
also to the varying integration intervals. The three limiting
integrals are
\[
 \int_{-1}^1(1-t^2)\dd t=\frac43,\qquad
 \int_{-1}^1(1-t^2)^2\dd t=\frac{16}{15},\qquad
 \int_{-1}^1t^2\dd t=\frac23.
\]
Restoring the factors \(\sqrt{2\epsilon}\), \(\epsilon\) and
\(\sqrt{2\epsilon}\) proves \eqref{eq:lobe-L1}--\eqref{eq:lobe-T}.
Their quotient gives \eqref{eq:lobe-Rayleigh}.
\end{proof}

The angular-integrated field mass and radial kinetic form in this
interval are \((2/\pi)M_{\mathrm{loc}}\) and
\((2/\pi)T_{\mathrm{loc}}\). Their ratio is unchanged. Both contributions
vanish at the collision. The divergence in \eqref{eq:lobe-Rayleigh}
is produced by normalizing a region whose mass vanishes faster than
its kinetic form; it is compatible with all the global regularity and
dissipation laws in Proposition~\ref{prop:global-energy}.

This quotient is a local Dirichlet quadratic form. The restriction of
\(h\) to \(I_\tau\), extended by zero, belongs to \(H^1(\R)\), since its
boundary values vanish. Its derivative generally jumps at the simple
endpoints, so the extension need not belong to \(H^2(\R)=\Dom B\).
Thus \eqref{eq:lobe-Rayleigh} is not being promoted to an operator
expectation of an arbitrarily cut-off state in the full domain.
Nor is a physical conditioning operation introduced by the
normalization.

At the event itself, \eqref{eq:heat-four-jet} gives
\begin{equation}\label{eq:local-density-scales}
 h(\tau_*,x_*+y)^2=\frac{a^2}{4}y^4+O(y^5),\qquad
 h(\tau_*+s,x_*)^2=a^2s^2+O(s^3).
\end{equation}
The local intensity therefore develops a quartic spatial zero and
refills quadratically in time, consistently with
\eqref{eq:density-event}. In the real amplitude, two sign transitions
disappear when the intervening pocket closes.

It is useful to compare three different notions of size near the same
event. The full norm and field energy are finite; the local lobe mass
and energy vanish; and the normalized local quotient diverges. Singular
zero-configuration quantities provide another example:
\[
 -\log\delta\sim-\tfrac12\log(8\epsilon),\qquad
 \delta^{-2}\sim(8\epsilon)^{-1}.
\]
These are functions of the zero coordinates, not the field expectation
\(\inner{\Psi}{K\Psi}\). The formal global zero Hamiltonians and
energies in \cite{RodgersTao2020} require their own truncation and
convergence analysis. Their singularity cannot be identified with a
failure of the smooth cylinder evolution.


\subsection{Exact balances on a moving sign region}
The asymptotic constants fit exact identities valid before the
collision. Write \(L_{\mathrm{loc}}=\int_{I_\tau}|h|\dd x\).
\begin{proposition}[Moving-boundary dissipation]\label{prop:lobe-balance}
As long as the endpoints are simple zeros bounding one sign region,
\begin{align}
 L_{\mathrm{loc}}'&=-|h_x(\tau,r_-)|-|h_x(\tau,r_+)|,
 \label{eq:lobe-L-balance}\\
 M_{\mathrm{loc}}'&=-2T_{\mathrm{loc}},
 \qquad \frac{T_{\mathrm{loc}}}{M_{\mathrm{loc}}}
       =-\tfrac12(\log M_{\mathrm{loc}})',
 \label{eq:lobe-M-balance}\\
 T_{\mathrm{loc}}'&=[h_xh_{xx}]_{r_-}^{r_+}
             -2\int_{I_\tau}h_{xx}^2\dd x.
 \label{eq:lobe-T-balance}
\end{align}
\end{proposition}
\begin{proof}
The boundary values of \(|h|\) and \(h^2\) vanish. If \(\sigma\)
is the sign of \(h\) in the region, differentiating its area gives
\(\sigma\int h_{xx}=\sigma[h_x]_{r_-}^{r_+}\), equal to
the negative sum in \eqref{eq:lobe-L-balance}.
Integration by parts gives
\(M_{\mathrm{loc}}'=2\int hh_{xx}=-2\int h_x^2\).
For the kinetic form, retain the endpoint velocities:
\[
 T_{\mathrm{loc}}'=h_x(r_+)^2r_+'-h_x(r_-)^2r_-'
                       +2\int_{I_\tau}h_xh_{xxx}\dd x.
\]
Using \(r_\pm'=-h_{xx}/h_x\) and integrating the last term
by parts leaves exactly \eqref{eq:lobe-T-balance}.
\end{proof}

The boundary term in \eqref{eq:lobe-T-balance} is essential: the
whole-line kinetic dissipation law does not apply unchanged to a
moving region. Equations \eqref{eq:lobe-M-balance} and
\eqref{eq:lobe-M} are consistent with the coefficient \(5/4\)
in \eqref{eq:lobe-Rayleigh}; that coefficient was proved from the
two integrals, without differentiating an uncontrolled asymptotic
remainder. In terms of the actual gap,
\begin{equation}\label{eq:lobe-gap-scales}
 L_{\mathrm{loc}}\sim\frac{|a|\delta^3}{12},\quad
 M_{\mathrm{loc}}\sim\frac{a^2\delta^5}{120},\quad
 T_{\mathrm{loc}}\sim\frac{a^2\delta^3}{12},\quad
 \delta^2\frac{T_{\mathrm{loc}}}{M_{\mathrm{loc}}}\longrightarrow10.
\end{equation}
These follow from \(\delta^2\sim8\epsilon\). Here
\(\epsilon=\tau_*-\tau=\lambda-\lambda_*\) is remaining time
to the event, not an error term or a source perturbation.
